\documentclass[a4paper,12pt]{article}
\usepackage[T2A]{fontenc}
\usepackage[utf8]{inputenc}
\usepackage[russian]{babel}
\usepackage{amsmath,amssymb,icomma,booktabs,tabularx,array,xcolor,tikz,fancyhdr}
\usepackage[a4paper,margin=23mm,headheight=15pt]{geometry}
\usepackage[colorlinks=true,linkcolor=teal!70!black]{hyperref}
\definecolor{accent}{HTML}{008080}
\definecolor{darkaccent}{HTML}{004D4D}
\definecolor{err}{HTML}{9A3737}
\setlength{\parindent}{0pt}
\setlength{\parskip}{5pt}
\pagestyle{fancy}
\fancyhf{}
\fancyhead[L]{\small Движение по окружности и тяготение}
\fancyhead[R]{\small София Лузянина}
\fancyfoot[C]{\thepage}
\renewcommand{\headrulewidth}{0.3pt}
\newcommand{\topic}[1]{\par\medskip\textbf{\textcolor{darkaccent}{#1}}\par}
\newcommand{\error}[1]{\par\medskip\textbf{\textcolor{err}{Ошибка:}} #1\par}
\newcommand{\key}[1]{\par\medskip\textbf{\textcolor{blue!60!black}{Ключевая идея:}} #1\par}
\newcommand{\result}[1]{\par\medskip\textbf{\textcolor{accent}{Итог:}} #1\par}
\newcommand{\task}[1]{\clearpage\phantomsection\label{task:#1}\section*{Задание №\,#1}}
\begin{document}
\begin{titlepage}
\centering
\vspace*{36mm}
{\Huge Ревью домашней работы\par}
\vspace{25mm}
{\Large Лёвин Артём Александрович\par}
\vspace{12mm}
{\large 27 ноября 2026 г.\par}
\vfill
{\small Лёвин Артём Александрович\par эксклюзивно для Софии Лузяниной\par}
\vspace{18mm}
\begin{tikzpicture}
\draw[accent,line width=.9pt] (0,0) .. controls (2,0.7) and (4,-0.7) .. (6,0) .. controls (8,0.7) and (10,-0.7) .. (12,0);
\end{tikzpicture}
\end{titlepage}
\setcounter{page}{1}
\section*{Сводная таблица}
Включены решения, в которых выявлены ошибки в формулах или математических переходах. Численные ответы в этих заданиях совпадают с эталонными.
\par\medskip\noindent
{\renewcommand{\arraystretch}{1.55}\small
\begin{tabularx}{\linewidth}{@{}c >{\raggedright\arraybackslash}p{.17\linewidth} >{\raggedright\arraybackslash}X >{\raggedright\arraybackslash}p{.19\linewidth} >{\raggedright\arraybackslash}p{.19\linewidth}@{}}
\toprule
№ & Тема & Суть ошибки & Ваш ответ & Правильный ответ\\
\midrule
\hyperref[task:5]{5} & Всемирное тяготение & Пропущен множитель $G$ в формуле силы. & Уменьшилась в 6 раз & Уменьшилась в 6 раз\\
\hyperref[task:9]{9} & Сравнение орбит & Потеря квадрата радиуса; неверное выражение скорости. & $v_1/v_2=3$; $\nu_1/\nu_2=27$; $T_1/T_2=1/27$ & $v_1/v_2=3$; $\nu_1/\nu_2=27$; $T_1/T_2=1/27$\\
\hyperref[task:10]{10} & Плотность и период & В записи отношения ускорений между множителями стоит знак равенства. & $R_B/R_A=2$; $T_B/T_A=0{,}5$ & $R_B/R_A=2$; $T_B/T_A=0{,}5$\\
\bottomrule
\end{tabularx}}
\clearpage
\thispagestyle{empty}
\null\vfill\begin{center}{\Huge Разбор ошибок}\end{center}\vfill\null
\task{5}
\topic{Условие}
Массы двух тел увеличили в 2 и 3 раза, расстояние между их центрами --- в 6 раз. Как изменилась сила притяжения?
\topic{Ваше решение и ответ}
Вы записали формулу $F=m_1m_2/R^2$ и получили
\[ F'=\frac{6m_1m_2}{36R^2}=\frac16F. \]
Ваш ответ: сила уменьшилась в 6 раз.
\error{В исходной формуле отсутствует гравитационная постоянная $G$.}
Последний правильный шаг: выделены коэффициенты изменения масс и расстояния. Первый неверный шаг --- запись абсолютной силы без $G$. В законе всемирного тяготения эта постоянная обязательна:
\[ F=G\frac{m_1m_2}{R^2}. \]
Без неё правая часть имеет размерность $\text{кг}^2/\text{м}^2$, которая отличается от размерности силы --- ньютона. Множитель $G$ обеспечивает правильную размерность и численное значение силы.
\key{При сравнении двух сил постоянная $G$ сокращается в их отношении.}
\topic{Исправление и полное решение}
Обозначим исходные массы и расстояние через $m_1,m_2,R$. После изменения они равны $2m_1,3m_2,6R$. Запишем новую силу:
\[ F'=G\frac{(2m_1)(3m_2)}{(6R)^2}=G\frac{6m_1m_2}{36R^2}. \]
Делим новую силу на исходную:
\[ \frac{F'}{F}=\frac{6Gm_1m_2}{36R^2}\cdot\frac{R^2}{Gm_1m_2}=\frac16. \]
\topic{Проверка}
Произведение масс возросло в 6 раз, квадрат расстояния --- в 36 раз. Общий коэффициент равен $6/36=1/6$.
\result{Сила уменьшилась в 6 раз. Ваш численный ответ верен; исправьте запись закона.}
\topic{Задача на закрепление}
Массы двух тел увеличили в 3 и 4 раза, а расстояние между их центрами --- в 6 раз. Определите, во сколько раз и в какую сторону изменилась сила притяжения.
\task{9}
\topic{Условие}
Два спутника вращаются вокруг одной планеты по круговым орбитам. Радиус второй орбиты в 9 раз больше радиуса первой. Во сколько раз отличаются линейные скорости, частоты и периоды обращения? Укажите, какая величина больше и у какого спутника.
\topic{Ваше решение и ответ}
Вы верно записали $R_2=9R_1$ и $a_{\text{ц}}=v^2/R$. Далее в законе тяготения использовали $R_1$ в первой степени и получили $a_1/a_2=9$. Затем записали $v^2=a_{\text{ц}}/R$ и при сравнении скоростей сократили радиусы. В конце получены отношения $v_1/v_2=3$, $\nu_1/\nu_2=27$, $T_1/T_2=1/27$.
\error{Две ошибки в формулах и необоснованное сокращение разных радиусов.}
Первый неверный шаг после записи условия --- равенство $ma_1=GMm/R_1$. Расстояние от центра планеты в законе тяготения возводится в квадрат. Поэтому
\[ ma_1=\frac{GMm}{R_1^2},\qquad a_1=\frac{GM}{R_1^2}. \]
Для второго спутника аналогично $a_2=GM/R_2^2$. Отсюда
\[ \frac{a_1}{a_2}=\frac{R_2^2}{R_1^2}=9^2=81. \]
Вторая ошибка возникает при выделении скорости. Из правильного равенства $a_{\text{ц}}=v^2/R$ следует после умножения обеих частей на $R$:
\[ v^2=a_{\text{ц}}R,\qquad v=\sqrt{a_{\text{ц}}R}. \]
Запись $v=\sqrt{a_{\text{ц}}/R}$ имеет размерность обратного времени и соответствует угловой скорости. Кроме того, радиусы двух орбит различны и сокращать их как одинаковые величины нельзя.
\key{Для круговой орбиты сила тяготения сообщает спутнику центростремительное ускорение. Сначала выведите формулы через радиус и затем сравнивайте величины.}
\topic{Исправление и полное решение}
Пусть $M$ --- масса планеты, $m$ --- масса спутника, $G$ --- гравитационная постоянная. По второму закону Ньютона:
\[ m\frac{v^2}{R}=G\frac{Mm}{R^2}. \]
Сокращаем массу спутника и умножаем обе части на $R$:
\[ v^2=\frac{GM}{R},\qquad v=\sqrt{\frac{GM}{R}}. \]
\clearpage
\subsection*{Задание №\,9. Продолжение}
\textbf{1. Линейные скорости.} Планета одна, поэтому $G$ и $M$ одинаковы для двух спутников:
\[ \frac{v_1}{v_2}=\sqrt{\frac{GM}{R_1}\cdot\frac{R_2}{GM}}=\sqrt{\frac{R_2}{R_1}}=\sqrt9=3. \]
Скорость первого спутника больше в 3 раза.

\textbf{2. Частоты обращения.} За один оборот спутник проходит путь $2\pi R$. Частота $\nu=v/(2\pi R)$. Поэтому
\[ \frac{\nu_1}{\nu_2}=\frac{v_1}{v_2}\cdot\frac{R_2}{R_1}=3\cdot9=27. \]
Частота первого спутника больше в 27 раз. Можно также использовать $\omega=2\pi\nu$: отношение частот равно отношению угловых скоростей.

\textbf{3. Периоды обращения.} Период и частота взаимно обратны:
\[ T=\frac1\nu,\qquad \frac{T_1}{T_2}=\frac{\nu_2}{\nu_1}=\frac1{27}. \]
Период второго спутника больше в 27 раз.
\topic{Проверка}
Независимо проверим отношение ускорений по формуле $a=v^2/R$:
\[ \frac{a_1}{a_2}=\left(\frac{v_1}{v_2}\right)^2\frac{R_2}{R_1}=3^2\cdot9=81. \]
Это совпадает с законом тяготения. Также $T_2/T_1=(R_2/R_1)/(v_2/v_1)=27$.
\result{Первый спутник движется в 3 раза быстрее и совершает в 27 раз больше оборотов за одинаковое время; период второго больше в 27 раз. Итоговые числа в Вашей работе совпали с верными, однако ошибочные промежуточные зависимости необходимо исправить.}
\topic{Задача на закрепление}
Два спутника движутся по круговым орбитам вокруг одной планеты. Радиус второй орбиты в 16 раз больше радиуса первой. Найдите отношения центростремительных ускорений, линейных скоростей, частот и периодов обращения первого спутника ко второму. Укажите, у какого спутника каждая величина больше.
\task{10}
\topic{Условие}
Средняя плотность сферической планеты B в 4 раза больше плотности планеты A, её масса в 32 раза больше. Найдите отношения $R_B/R_A$ и $T_B/T_A$ для низких круговых орбит, радиусы которых равны радиусам планет.
\topic{Ваше решение и ответ}
Из $V=M/\rho$ и $V=4\pi R^3/3$ Вы верно получили $V_B/V_A=8$ и $R_B=2R_A$. Далее использовали $a=GM/R^2$ и получили $a_B/a_A=8$, затем $\omega_B/\omega_A=2$ и $T_B/T_A=0{,}5$.
\error{При записи отношения ускорений между дробями поставлен знак равенства вместо знака умножения.}
Последний корректный шаг --- $R_B=2R_A$ и формула $a=GM/R^2$. В строке с отношением ускорений две дроби соединены знаком равенства; исправление коэффициента в знаменателе частично затёрто. По дальнейшему незачёркнутому результату видно, что коэффициент 8 получен верно. Ошибку в самом коэффициенте не утверждаем.
\key{При делении на дробь необходимо умножить на обратную ей дробь. Квадрат относится ко всему радиусу: $(2R_A)^2=4R_A^2$.}
\topic{Исправление и полное решение}
Из условия $M_B=32M_A$, $\rho_B=4\rho_A$. Тогда
\[ \frac{V_B}{V_A}=\frac{M_B/\rho_B}{M_A/\rho_A}=\frac{32}{4}=8. \]
Объём шара пропорционален кубу радиуса, следовательно
\[ \left(\frac{R_B}{R_A}\right)^3=8,\qquad \frac{R_B}{R_A}=2. \]
На низких круговых орбитах $a=GM/R^2$. Запишем деление полностью:
\[ \frac{a_B}{a_A}=\frac{G\cdot32M_A}{(2R_A)^2}\cdot\frac{R_A^2}{GM_A}=\frac{32}{4}=8. \]
Из $a=\omega^2 R$ находим $\omega=\sqrt{a/R}$. Тогда
\[ \frac{\omega_B}{\omega_A}=\sqrt{\frac{a_B}{a_A}\cdot\frac{R_A}{R_B}}=\sqrt{8\cdot\frac12}=2. \]
\clearpage
\subsection*{Задание №\,10. Продолжение}
Поскольку $T=2\pi/\omega$, получаем
\[ \frac{T_B}{T_A}=\frac{\omega_A}{\omega_B}=\frac12=0{,}5. \]
\topic{Проверка другим способом}
Из равенства $\omega^2R=GM/R^2$ следует
\[ T=2\pi\sqrt{\frac{R^3}{GM}}. \]
Для сферической планеты $M=\rho V=4\pi\rho R^3/3$. Подставим массу в формулу периода:
\[ T=2\pi\sqrt{\frac{3}{4\pi G\rho}}=\sqrt{\frac{3\pi}{G\rho}}. \]
Следовательно, при принятых в условии низких круговых орбитах период зависит от средней плотности планеты:
\[ \frac{T_B}{T_A}=\sqrt{\frac{\rho_A}{\rho_B}}=\sqrt{\frac14}=\frac12. \]
Результат совпадает с первым способом. Это также объясняет, почему увеличение плотности планеты в четыре раза уменьшает период вдвое при орбите у поверхности.
\result{Радиус планеты B в 2 раза больше; период спутника у её поверхности в 2 раза меньше. Оба Ваших ответа верны. Исправьте знак операции в записи отношения ускорений.}
\topic{Задача на закрепление}
Средняя плотность сферической планеты B в 9 раз больше плотности планеты A, а её масса в 72 раза больше. Найдите $R_B/R_A$ и $T_B/T_A$ для спутников на низких круговых орбитах, радиусы которых равны радиусам соответствующих планет.
\vfill
\begin{center}\small Лёвин Артём Александрович\par эксклюзивно для Софии Лузяниной\end{center}
\end{document}
